\chapter{Conclusions and Outlook}

\section{Conclusions}

The aim of this thesis has been to investigate how effective two-flavor
quark-matter models can be used to construct compact-star equations of state and
study the resulting mass--radius relations. Building on the compact-star
framework developed in the project thesis, the focus here was shifted from
hadronic matter to deconfined quark degrees of freedom. The main question was
whether chiral quark models, and in particular the QMD model with two-flavor
color-superconducting pairing, can produce compact-star
sequences compatible with the main qualitative astrophysical constraints.

The two-flavor Quark--Meson model was first used as a baseline. After imposing a
vacuum-stability check, the model was studied for
\(m_\sigma=400,500,\) and \(600~\mathrm{MeV}\). Lower values of \(m_\sigma\)
produce stiffer equations of state and larger maximum masses, while the
\(m_\sigma=400~\mathrm{MeV}\) and \(m_\sigma=500~\mathrm{MeV}\) cases display
compact, self-bound-like low-mass behavior associated with a first-order
transition and finite surface density. By contrast, the
\(m_\sigma=600~\mathrm{MeV}\) case gives a smoother low-density behavior and a
gravitationally bound branch. This shows that even in the minimal two-flavor
Quark--Meson model, the chiral sector affects not only the mass and radius
scale, but also the topology of the stellar sequence.

A Bodmer--Witten-inspired shift was then introduced as a phenomenological
bag-like modification of the vacuum normalization. Since the model contains only
up and down quarks, this should not be interpreted as a calculation of
absolutely stable strange matter. Its role in this thesis was instead to test
how the quark-star sequences respond to changes in the zero-pressure point. The
shift reduces the characteristic mass and radius scales and makes some
sequences less extreme, but it does not qualitatively change the ordering or
topology in the explored range: the \(m_\sigma=400~\mathrm{MeV}\) and
\(m_\sigma=500~\mathrm{MeV}\) branches remain self-bound-like, while the
\(m_\sigma=600~\mathrm{MeV}\) branch remains gravitationally bound. These
results are useful diagnostics, but they mainly serve as a stepping stone toward
the color-superconducting QMD model.

The central result of the thesis comes from the two-flavor QMD model. By
introducing a diquark condensate, the model
allows two-flavor color-superconducting matter to be studied explicitly. A key
technical conclusion is that the full one-loop potential is necessary for a
robust 2SC phase. In particular, the finite residual contribution
\(\Omega_{1,\mathrm{num}}\) cannot be treated as a harmless numerical correction:
removing it qualitatively changes the paired phase structure. In the truncated
potential, the 2SC phase may collapse into a narrow finite window or disappear
entirely, depending on the parameter choice. All final QMD results therefore
have to be based on the full one-loop potential.

With the full potential, the QMD model develops a robust 2SC phase both in the
common-\(\mu_q\) benchmark and in neutral stellar matter. In the stellar case,
the onset of pairing is shifted relative to the common-\(\mu_q\) calculation by
the neutrality conditions, but the qualitative effect remains the same: pairing
stiffens the equation of state in the density range relevant for compact stars.
The baseline neutral QMD sequence is gravitationally bound and reaches
\[
    M_{\max}=1.970~M_\odot,
    \qquad
    R(M_{\max})=12.16~\mathrm{km}.
\]
This places the baseline model close to the two-solar-mass scale and shows that
2SC pairing can have a sizeable effect on the stellar structure of two-flavor
quark matter.

The parameter study showed that the QMD model contains several complementary
directions for changing the mass--radius relation. Increasing the diquark Yukawa
coupling \(g_\Delta\) raises the maximum mass in the stellar density range,
while lowering the diquark mass parameter \(m_\Delta\) shifts the 2SC phase into
the relevant density range earlier and also increases the maximum mass. The
self-coupling \(\lambda_\Delta\) has only a small effect in the explored range.
The mixed coupling \(\lambda_3\), on the other hand, mainly controls the radius
scale and the shape of the stellar branch: increasing it compresses the sequence
and can improve agreement with smaller-radius observational regions, although
too large a value can lead to less clean branch behavior.

The most promising phenomenology was obtained by combining parameters that act
in complementary ways. In particular, one can combine parameters that raise the
maximum mass with parameters that reduce the radius scale. Two such combined
models exceed \(2M_\odot\) while giving radii of roughly \(11.8\)--\(11.9\) km
at the maximum mass, and both pass visually through the plotted NICER regions
for J0030+0451 and J0740+6620. A third combined model also exceeds
\(2M_\odot\), but becomes too compact for J0740+6620 and develops a less clean,
partly multivalued mass--radius topology. The combined tests therefore indicate
that the QMD model can produce simultaneously massive and moderately compact
stellar sequences, but only in a restricted region of parameter space.

The comparison with astrophysical constraints should still be interpreted
cautiously. It is visual only and does not include tidal deformabilities,
Bayesian inference, or a systematic multidimensional parameter scan. The models
studied here are also pure quark-star sequences. They do not contain a hadronic
crust, a nuclear envelope, or a transition from hadronic to quark matter.
Furthermore, the model contains only two quark flavors and therefore omits
strange quarks, CFL pairing, and other degrees of freedom expected to be
important in a more complete description of dense matter. None of the explored
models reach the highest-mass candidate bands considered.

The main conclusion of the thesis is therefore that two-flavor QMD matter with
2SC pairing can produce phenomenologically promising compact-star sequences,
provided the full one-loop potential is retained. The baseline model already
lies close to the two-solar-mass scale, and selected combined parameter choices
exceed \(2M_\odot\) while giving radii compatible with the plotted NICER regions
for both J0030+0451 and J0740+6620. At the same time, these results should not
be interpreted as a final realistic neutron-star equation of state. Rather, they
identify color-superconducting QMD matter as a promising component for future
hybrid-star and three-flavor compact-star modelling.

\section{Outlook}

The most direct continuation of this thesis is to move beyond pure quark-star
sequences and embed the QMD phase in a hybrid-star construction. In the present
work, the stellar configurations are composed entirely of quark matter, which is
useful for isolating the role of the quark-sector parameters, but not sufficient
for a realistic description of observed compact stars. A hybrid construction
would allow the low- and intermediate-density regions to be described by nuclear
matter, while the QMD equation of state would enter above a transition density
where deconfined quark degrees of freedom become relevant
\cite{Baym2018,Alford2013,Annala2020,Andersen2026}. This would also make it
possible to include a crust or nuclear envelope, which is absent in the present
pure quark-star treatment.

A first implementation could use a Maxwell construction between a hadronic
equation of state and the QMD equation of state. The resulting hybrid-star
sequences could then be tested against the generic stability criteria for compact
stars with sharp phase transitions \cite{Alford2013}. More advanced treatments
could consider Gibbs constructions or crossover interpolations, depending on the
assumed surface tension and microscopic nature of the hadron--quark transition
\cite{Baym2018,Masuda2013,Kurkela2014}. This would determine whether the
promising QMD parameter combinations found in this thesis produce stable hybrid
branches, disconnected twin-star branches, or only small quark cores when
embedded in a hadronic background.

A second important direction is to extend the model from two to three quark
flavors. The present work includes only up and down quarks, so the
Bodmer--Witten-inspired shift should not be interpreted as a calculation of
absolutely stable strange matter. A genuine discussion of strange quark matter
requires the strange quark and the associated condensates to be included
dynamically \cite{Bodmer1971,Witten1984,Baym2018}. This would also make it
possible to study the color-flavor-locked (CFL) phase, which is expected to be favoured
in three-flavor quark matter at asymptotically large densities
\cite{Alford2008}. A three-flavor QMD model would therefore test whether the
promising two-flavor parameter directions found here survive when strange quarks
and competition between 2SC and CFL pairing are included
\cite{Alford2008,Andersen2026}.

Another natural extension is the inclusion of vector and axial-vector mesons.
Repulsive vector interactions can increase the stiffness of dense quark matter
and are therefore relevant for the compatibility of quark matter with massive
compact stars \cite{Klahn2015,Otto2020}. In the present thesis, stiffness is
generated mainly through the chiral and diquark sectors. Adding vector degrees
of freedom would provide an additional mechanism for increasing the pressure at
high density and would help disentangle the effects of diquark pairing from
repulsive quark interactions. It would also make it possible to study the speed
of sound more systematically, especially since massive compact stars may require
the speed of sound to exceed the conformal value over some density interval
\cite{Bedaque2015,Tews2018,Annala2020,Komoltsev2022}.

The comparison with observations should also be made more quantitative. In this
thesis, the comparison with massive pulsars and NICER regions was visual and
diagnostic. A natural next step is to compute likelihoods or posterior
distributions using mass measurements, NICER pulse-profile constraints, and
gravitational-wave tidal deformability constraints
\cite{Riley_2019_J0030,Miller_2019_J0030,Miller_2021_J0740,Riley2021,
LIGO2018}. Tidal deformabilities are especially important because they are
highly sensitive to the stellar radius and therefore provide information
complementary to the maximum mass \cite{LIGO2018,De2018,Raithel2018,Tews2018}.
This would test whether the combined QMD parameter sets that visually overlap
the NICER regions are also compatible with binary-neutron-star merger data.

A systematic multidimensional parameter scan would be particularly useful. The
results of this thesis show that \(g_\Delta\), \(m_\Delta\), and \(\lambda_3\)
affect complementary aspects of the mass--radius relation: some directions raise
the maximum mass, while others reduce the radius scale or change the branch
shape. The successful combined parameter sets should therefore be interpreted as
promising points in parameter space, not as optimized models. A broader scan
around these regions, with constraints from stability, causality, chiral
restoration, and high-density QCD, would clarify whether the viable region is
broad or fine-tuned \cite{Kurkela2010,Kurkela2014,Komoltsev2022,Annala2020}.

Finally, the role of \(\Omega_{1,\mathrm{num}}\) deserves further study. In this
thesis, this term was shown to be structurally necessary for robust 2SC pairing.
A more detailed analysis of its density dependence and its effect on the
curvature of the thermodynamic potential could clarify why it has such a large
impact on the paired phase. This would strengthen the theoretical understanding
of the QMD model and help determine whether similar residual one-loop
contributions are important in related effective models of dense quark matter.

In summary, the most important next steps are a hybrid-star construction using
the full one-loop QMD potential, a three-flavor extension including strange
quarks and CFL pairing, and a systematic parameter scan constrained by masses,
radii, and tidal deformabilities. Together, these developments would turn the
present two-flavor pure-quark-star study into a more realistic framework for
testing whether color-superconducting quark matter can exist inside observed
compact stars. In this direction, the more complete hybrid and three-flavor QMD
construction of Ref.~\cite{Andersen2026} provides a natural point of comparison
and a concrete example of how several of these extensions can be implemented.
